% !TEX program = xelatex
\documentclass[UTF8,10pt]{ctexart}
\usepackage[margin=1in]{geometry}
\usepackage{amsmath,amssymb}
\usepackage{booktabs}
\usepackage{tabularx}
\usepackage{graphicx}
\usepackage{float}
\usepackage{caption}
\usepackage{hyperref}
\usepackage{xcolor}
\usepackage{lmodern}
\captionsetup{font=small,labelfont=bf}
\setlength{\parindent}{0pt}
\setlength{\parskip}{0.5em}
\newcommand{\best}[1]{\textbf{#1}}
\newcommand{\second}[1]{\underline{#1}}
\begin{document}
\title{视频风格迁移评估 - 合成流程演示}
\author{自动生成评估报告}
\date{2026-11-08 \quad \textbf{结论状态: 探索性}}
\maketitle
\begin{abstract}
演示可复现的风格、内容与运动评估流程，并明确合成片段不能作为模型基准证据。
\par
本报告用 3 个合成案例验证了评估流程，未提供真实模型性能证据。结果表、感知分析和结论仍保留多维证据，不应被解释为 EchoStyle 或本文方法的基准排名。
\end{abstract}
\section*{评估设置与可复现性}
\textbf{证据完整性检查}
\begin{center}
\small
\begin{tabular}{@{}lll@{}}
\toprule
检查项 & 结果 & 状态 \\
\midrule
有效配对案例 & 3 & 通过 \\
方法数量 & 3 & 通过 \\
盲评 VLM 记录 & 已提供 & 通过 \\
论文级感知指标 & 未检测到 & 尚未建立 \\
结论状态 & 探索性 & 需说明限制 \\
\bottomrule
\end{tabular}
\end{center}
\begin{itemize}
\item 光流实现: OpenCV Farneback (diagnostic proxy)
\item 每个视频采样帧数: 30
\item 时序偏移: [1, 5]
\item 盲评: True
\end{itemize}
\section*{指标定义与设计动机}
\subsection*{SSIM 结构相似度}
\textit{在本次运行中计算}
\begin{equation*}
\mathrm{SSIM}(x,y)=\frac{(2\mu_x\mu_y+C_1)(2\sigma_{xy}+C_2)}{(\mu_x^2+\mu_y^2+C_1)(\sigma_x^2+\sigma_y^2+C_2)}
\end{equation*}
\textbf{设计动机.} SSIM 比较局部亮度、对比度和结构，而不是原始像素误差。它适合检测几何和结构变化，但强风格化本身也会改变对比度与纹理。
\textbf{解释与报告注意事项.} 应将其视为内容诊断，而不是感知风格分数。
\subsection*{PSNR 峰值信噪比}
\textit{在本次运行中计算}
\begin{equation*}
\mathrm{PSNR}=10\log_{10}\!\left(\frac{\mathrm{MAX}^2}{\mathrm{MSE}}\right)
\end{equation*}
\textbf{设计动机.} PSNR 将均方误差转换为对数尺度，提供快速、确定性的保真度诊断。它对全局颜色变化、压缩和分辨率差异较敏感。
\textbf{解释与报告注意事项.} 不能把 PSNR 的小幅变化直接解释为感知质量的等比例变化。
\subsection*{LPIPS 感知距离}
\textit{本次运行未提供；定义用于论文级扩展}
\begin{equation*}
\mathrm{LPIPS}(x,y)=\sum_l\frac{1}{H_lW_l}\sum_{h,w}\left\|w_l\odot(\hat\phi_l(x)_{h,w}-\hat\phi_l(y)_{h,w})\right\|_2^2
\end{equation*}
\textbf{设计动机.} LPIPS 比较归一化深层特征，通常比像素指标更符合人类感知相似性。在有固定 LPIPS 实现时，应优先作为内容保持证据。
\textbf{解释与报告注意事项.} 所有比较方法必须使用相同的网络、预处理和权重。
\subsection*{CLIP-F 内容语义相似度}
\textit{本次运行未提供；定义用于论文级扩展}
\begin{equation*}
\mathrm{CLIP\text{-}F}(R,I)=100\cos\!\left(\phi_I(R),\phi_I(I)\right)
\end{equation*}
\textbf{设计动机.} CLIP-F 衡量风格化帧与源帧是否保留相同的图像级语义。当外观变化较大时，它可以补充结构指标。
\textbf{解释与报告注意事项.} 图像级嵌入可能遗漏局部身份、边界和遮挡错误。
\subsection*{CLIP-I 风格图像相似度}
\textit{本次运行未提供；定义用于论文级扩展}
\begin{equation*}
\mathrm{CLIP\text{-}I}(R,S)=100\cos\!\left(\phi_I(R),\phi_I(S)\right)
\end{equation*}
\textbf{设计动机.} CLIP-I 提供对内容变化较稳健的全局信号，用于判断生成帧与风格参考是否位于相近的语义视觉空间。它只是风格证据的一部分，不能单独证明纹理迁移成功。
\textbf{解释与报告注意事项.} 应进行相对比较，并结合纹理特征或盲评审查。
\subsection*{FID 分布距离}
\textit{本次运行未提供；定义用于论文级扩展}
\begin{equation*}
\mathrm{FID}=\left\|\mu_s-\mu_g\right\|_2^2+\operatorname{Tr}\!\left(\Sigma_s+\Sigma_g-2(\Sigma_s\Sigma_g)^{1/2}\right)
\end{equation*}
\textbf{设计动机.} FID 比较 Inception 特征分布，因此能捕捉超出单图颜色统计的风格集合差异。它需要足够数量的参考和生成帧来估计分布。
\textbf{解释与报告注意事项.} 单张风格图不能构成有效的 FID 风格分布。必须记录图像数量、分辨率和预处理。
\subsection*{ArtFID 联合风格-内容距离}
\textit{本次运行未提供；定义用于论文级扩展}
\begin{equation*}
\mathrm{ArtFID}=\left(1+\frac{1}{N}\sum_i\mathrm{LPIPS}(I_i,R_i)\right)\left(1+\mathrm{FID}(R,S)\right)
\end{equation*}
\textbf{设计动机.} ArtFID 将感知内容保持和风格分布距离结合为一个数值。乘积形式会让“一个因素改善、另一个因素退化”的方法暴露出来。
\textbf{解释与报告注意事项.} 除非 LPIPS 和 FID 分量遵循固定的引用实现，否则不要使用该指标。
\subsection*{LAB 调色板 W1}
\textit{在本次运行中计算}
\begin{equation*}
W_{1,c}=\frac{1}{255}\sum_{k=0}^{255}\left|F^{\mathrm{frame}}_c(k)-F^{\mathrm{style}}_c(k)\right|,\qquad \mathrm{PaletteW1}=100\cdot\frac{1}{|C|}\sum_{c\in C}W_{1,c}
\end{equation*}
\textbf{设计动机.} 该低成本的诊断指标概括 LAB 颜色边缘分布的一致性，可在运行昂贵感知指标前发现全局着色和调色板不匹配。
\textbf{解释与报告注意事项.} 它不衡量笔触、纹理附着、构图或风格语义。
\subsection*{源光流重映射 RMSE}
\textit{在本次运行中计算}
\begin{equation*}
\mathrm{WarpRMSE}_{\Delta}=\sqrt{\frac{\sum_x M_t(x)\left\|R_t(x)-W(R_{t+\Delta})(x)\right\|_2^2}{\sum_x M_t(x)}}
\end{equation*}
\textbf{设计动机.} 源视频光流提供预期运动。使用该光流将后续输出帧重映射回当前帧，可以将漂移和闪烁与正常场景运动分开，并屏蔽遮挡像素。
\textbf{解释与报告注意事项.} 结果会继承光流估计器的误差；Farneback 是诊断代理，不是 RAFT。
\subsection*{光流端点误差}
\textit{在本次运行中计算}
\begin{equation*}
\mathrm{EPE}_{\Delta}=\frac{1}{|\Omega|}\sum_{x\in\Omega}\left\|F^R_t(x)-F^I_t(x)\right\|_2
\end{equation*}
\textbf{设计动机.} EPE 衡量风格化是否改变帧间估计位移，可揭示静态图像难以发现的运动改变和传播失败。
\textbf{解释与报告注意事项.} 纹理变化可能产生虚假光流；所有方法必须使用相同估计器和遮挡处理。
\subsection*{光流场相关性}
\textit{在本次运行中计算}
\begin{equation*}
\rho_{\Delta}=\operatorname{corr}\!\left(\operatorname{vec}(F^R_t),\operatorname{vec}(F^I_t)\right)
\end{equation*}
\textbf{设计动机.} 相关性衡量运动方向的一致性，相比 EPE 更少受幅度影响。两者结合可以区分方向错误和位移尺度变化。
\textbf{解释与报告注意事项.} 这是一个全局概括，可能掩盖局部运动错误。
\subsection*{时序 LPIPS（tLP）}
\textit{本次运行未提供；定义用于论文级扩展}
\begin{equation*}
\mathrm{tLP}=\frac{1}{T-\Delta}\sum_t\mathrm{LPIPS}\!\left(R_t,W(R_{t+\Delta})\right)
\end{equation*}
\textbf{设计动机.} tLP 在重映射分析中用深层特征感知误差替代像素误差，使纹理和结构不稳定性比原始 L2 差值更容易被检测。
\textbf{解释与报告注意事项.} 必须使用固定的 LPIPS 网络并报告时间偏移。
\subsection*{CLIP 时序一致性}
\textit{本次运行未提供；定义用于论文级扩展}
\begin{equation*}
\mathrm{CLIP\text{-}CONS}=100\cos\!\left(\phi(R_{t+1})-\phi(R_t),\phi(I_{t+1})-\phi(I_t)\right)
\end{equation*}
\textbf{设计动机.} CLIP-CONS 检查相邻风格化帧的语义变化是否与源视频一致。当内容被强风格化时，它可以补充像素重映射指标。
\textbf{解释与报告注意事项.} 它衡量语义变化一致性，而不是空间对齐或闪烁。
\subsection*{去趋势亮度闪烁}
\textit{在本次运行中计算}
\begin{equation*}
F=\sqrt{\frac{1}{T}\sum_t\left(l_t-(at+b)\right)^2},\quad l_t=\operatorname{mean}(\operatorname{luma}(R_t))
\end{equation*}
\textbf{设计动机.} 去除线性亮度趋势可以区分合理的曝光变化和突发全局闪烁。该指标成本低，直接针对常见的视频风格迁移伪影。
\textbf{解释与报告注意事项.} 它忽略局部闪烁和纹理粘连，这些需要运动感知指标或视觉审查。
\subsection*{盲评 VLM 分数}
\textit{在本次运行中计算}
\begin{equation*}
\bar S_{d,m}=\frac{1}{R}\sum_{r=1}^{R}\frac{1}{|C_{\mathrm{valid}}|}\sum_{c\in C_{\mathrm{valid}}}s_{r,c,m,d},\quad s\in\{1,\ldots,5\}
\end{equation*}
\textbf{设计动机.} 该评分标准覆盖风格一致性、内容保持、时序稳定性、伪影控制和整体质量，能补充自动指标难以区分的纹理与运动故障。逐案例盲化方法标签可降低方法身份偏差。
\textbf{解释与报告注意事项.} VLM 判断是感知证据而非真实标签；应记录模型版本，并使用多个独立审查者支持论文级结论。
\subsection*{运行时间与峰值显存}
\textit{在本次运行中计算}
\begin{equation*}
\mathrm{Runtime}_m=\frac{1}{|C|}\sum_c t_{c,m},\qquad \mathrm{PeakMem}_m=\max_c p_{c,m}
\end{equation*}
\textbf{设计动机.} 如果缺少部署成本，质量比较并不完整。平均运行时间反映预期延迟，最大峰值显存反映最差观测资源需求。
\textbf{解释与报告注意事项.} 必须在相同硬件、精度、分辨率和批设置下测量，并同时报告均值和峰值。
\section*{数据与配对设计}
\begin{table}[H]
\centering
\caption{配对评估案例及有效方法子集。}
\label{tab:dataset}
\small
\begin{tabularx}{\linewidth}{@{}lXlX@{}}
\toprule
案例 & 内容 / 风格 & 数据划分 & 有效方法 \\
\midrule
demo\_case\_01 & content\_01 / style\_01 & synthetic\_demo & 模拟平衡方法, 模拟闪烁方法, 模拟覆盖方法 \\
demo\_case\_02 & content\_02 / style\_02 & synthetic\_demo & 模拟平衡方法, 模拟闪烁方法, 模拟覆盖方法 \\
demo\_case\_03 & content\_03 / style\_03 & synthetic\_demo & 模拟平衡方法, 模拟闪烁方法, 模拟覆盖方法 \\
\bottomrule
\end{tabularx}
\end{table}
\section*{实验结果}
\begin{table}[H]
\centering
\caption{风格迁移结果。表中为有效配对片段上的均值和样本标准差。最优和次优均值分别以粗体和下划线标出。箭头表示指标偏好方向。有效片段数：balanced: 3; flicker: 3; overlay: 3。}
\label{tab:metrics-1}
\small
\resizebox{\linewidth}{!}{%
\begin{tabular}{@{}lccc@{}}
\toprule
指标 & 模拟平衡方法 & 模拟闪烁方法 & 模拟覆盖方法 \\
\midrule
LAB 调色板 W1 $\downarrow$ & $9.0208 \pm 0.5445$ & \textbf{$7.4201 \pm 0.5347$} & \underline{$8.0416 \pm 0.9411$} \\
LAB 色度 W1 $\downarrow$ & \underline{$3.1765 \pm 1.3380$} & \textbf{$1.8824 \pm 1.0178$} & $4.6879 \pm 0.4238$ \\
\bottomrule
\end{tabular}%
}
\end{table}
\begin{table}[H]
\centering
\caption{内容保持结果。表中为有效配对片段上的均值和样本标准差。最优和次优均值分别以粗体和下划线标出。箭头表示指标偏好方向。有效片段数：balanced: 3; flicker: 3; overlay: 3。}
\label{tab:metrics-2}
\small
\resizebox{\linewidth}{!}{%
\begin{tabular}{@{}lccc@{}}
\toprule
指标 & 模拟平衡方法 & 模拟闪烁方法 & 模拟覆盖方法 \\
\midrule
SSIM 结构相似度 $\uparrow$ & \textbf{$0.9290 \pm 0.0165$} & \underline{$0.8111 \pm 0.0326$} & $0.7153 \pm 0.0149$ \\
PSNR 峰值信噪比 $\uparrow$ & \textbf{$22.58 \pm 0.9018$} & \underline{$17.54 \pm 0.3337$} & $14.90 \pm 0.3931$ \\
\bottomrule
\end{tabular}%
}
\end{table}
\begin{table}[H]
\centering
\caption{时序与运动结果。表中为有效配对片段上的均值和样本标准差。最优和次优均值分别以粗体和下划线标出。箭头表示指标偏好方向。有效片段数：balanced: 3; flicker: 3; overlay: 3。}
\label{tab:metrics-3}
\small
\resizebox{\linewidth}{!}{%
\begin{tabular}{@{}lccc@{}}
\toprule
指标 & 模拟平衡方法 & 模拟闪烁方法 & 模拟覆盖方法 \\
\midrule
重映射 RMSE Δ1 $\downarrow$ & \textbf{$31.70 \pm 6.7038$} & \underline{$47.77 \pm 10.70$} & $63.87 \pm 6.8744$ \\
重映射 RMSE Δ5 $\downarrow$ & \textbf{$33.20 \pm 5.2688$} & \underline{$46.18 \pm 10.55$} & $63.70 \pm 7.5960$ \\
光流端点误差 Δ1 $\downarrow$ & \underline{$0.2438 \pm 0.0558$} & $0.5426 \pm 0.3239$ & \textbf{$0.1874 \pm 0.0547$} \\
光流相关性 Δ1 $\uparrow$ & \textbf{$0.4268 \pm 0.2381$} & $0.2267 \pm 0.0923$ & \underline{$0.2662 \pm 0.1339$} \\
亮度闪烁 $\downarrow$ & \underline{$0.1035 \pm 0.0619$} & $12.43 \pm 0.4689$ & \textbf{$0.0822 \pm 0.0217$} \\
\bottomrule
\end{tabular}%
}
\end{table}
\begin{table}[H]
\centering
\caption{计算成本结果。表中为有效配对片段上的均值和样本标准差。最优和次优均值分别以粗体和下划线标出。箭头表示指标偏好方向。有效片段数：balanced: 3; flicker: 3; overlay: 3。}
\label{tab:metrics-4}
\small
\resizebox{\linewidth}{!}{%
\begin{tabular}{@{}lccc@{}}
\toprule
指标 & 模拟平衡方法 & 模拟闪烁方法 & 模拟覆盖方法 \\
\midrule
运行时间 $\downarrow$ & $18.80 \pm 0.4000$ & \textbf{$15.00 \pm 0.4000$} & \underline{$16.50 \pm 0.4000$} \\
峰值 GPU 显存 $\downarrow$ & $12.70 \pm 0.0000$ & \textbf{$9.8000 \pm 0.0000$} & \underline{$10.40 \pm 0.0000$} \\
\bottomrule
\end{tabular}%
}
\end{table}
\begin{table}[H]
\centering
\caption{相对主基线的配对比较。差值为方向校正后的结果，正值支持主方法。dz 为描述性配对标准化效应量；当 n 小于五时不报告置信区间。}
\label{tab:paired-effects}
\small
\begin{tabular}{@{}lrrrr@{}}
\toprule
评价指标 & $\Delta$ & $d_z$ & 95\% CI & $n$ \\
\midrule
SSIM 结构相似度 & 0.1180 & 7.1536 & -- & 3 \\
LAB 调色板 W1 & -1.6007 & -15.05 & -- & 3 \\
重映射 RMSE Δ1 & 16.08 & 3.9236 & -- & 3 \\
光流端点误差 Δ1 & 0.2988 & 1.0159 & -- & 3 \\
运行时间 & -3.8000 & n/a & -- & 3 \\
\bottomrule
\end{tabular}
\end{table}
\section*{盲评视觉分析}
\begin{table}[H]
\centering
\caption{盲评视觉评分（1--5），所有维度均为越高越好。}
\label{tab:vlm}
\small
\resizebox{\linewidth}{!}{%
\begin{tabular}{@{}lccccc@{}}
\toprule
方法 & 风格一致性 & 内容保持 & 时序稳定性 & 伪影控制 & 整体视觉质量 \\
\midrule
模拟平衡方法 & 4.00 & 4.67 & 3.67 & 3.33 & 4.00 \\
模拟闪烁方法 & 4.67 & 3.67 & 2.00 & 2.00 & 2.33 \\
模拟覆盖方法 & 3.33 & 3.00 & 3.00 & 2.00 & 2.00 \\
\bottomrule
\end{tabular}%
}
\end{table}
\noindent\textbf{失败标签:} 细节损失 (3), 纹理附着/粘连 (3), 闪烁 (3), 边缘重影 (2), 时序跳变 (2), 纹理漂移 (1), 几何融化 (1)\par
\section*{分析与讨论}
\subsection*{风格保真度}
现有指标总体上更支持 模拟闪烁方法。 已报告的诊断值为：LAB 调色板 W1 9.0208 ± 0.5445；LAB 色度 W1 3.1765 ± 1.3380。 盲评风格一致性为 4.00/5。
\subsection*{内容保持}
现有指标总体上更支持 模拟平衡方法。 已报告的诊断值为：SSIM 结构相似度 0.9290 ± 0.0165；PSNR 峰值信噪比 22.58 ± 0.9018。
\subsection*{运动与时序稳定性}
现有指标总体上更支持 模拟平衡方法。 已报告的诊断值为：重映射 RMSE Δ1 31.70 ± 6.7038；重映射 RMSE Δ5 33.20 ± 5.2688；光流端点误差 Δ1 0.2438 ± 0.0558；亮度闪烁 0.1035 ± 0.0619。
\subsection*{感知审查与失败模式}
盲评中整体视觉质量最高的是 模拟平衡方法（4.00/5，1 位审查者）。该结果仅代表当前审查样本中的感知偏好，不等同于真实标签。 主要失败模式为细节损失（3）、纹理附着/粘连（3）、闪烁（3）。
\subsection*{计算成本}
现有指标总体上更支持 模拟闪烁方法。 主方法的运行时间 18.80 ± 0.4000 秒、峰值显存 12.70 ± 0.0000 GB。
\subsection*{综合判断}
结果呈现明显的多维权衡，没有方法在所有维度上占优；因此不能据此声称单一方法的普遍优越性。
\section*{局限性与有效性威胁}
\begin{itemize}
\item 内容与风格的数据分布可能不能代表实际使用场景。
\item 低层指标容易受全局着色和压缩影响。
\item Farneback 光流不是 RAFT，在大运动或遮挡条件下可能失效。
\item 单张风格图不能构成有效的 FID 风格分布。
\item VLM 审查可能受到位置、内容和提示词偏差影响。
\item 直接引用论文的外部基线数值可能没有统一推理设置。
\item 少于五个片段时，比较置信区间不稳定。
\end{itemize}
\section*{结论}
本报告用 3 个合成案例验证了评估流程，未提供真实模型性能证据。结果表、感知分析和结论仍保留多维证据，不应被解释为 EchoStyle 或本文方法的基准排名。
\subsection*{建议}
请使用当前结果验证评估流程并设计下一组实验。不要把这个合成演示作为模型性能证据。
\clearpage
\section*{附录：单案例证据}
\subsection*{demo\_case\_01}
\begin{table}[H]
\centering
\caption{demo\_case\_01 的代表性指标。}
\small
\begin{tabular}{@{}lcccc@{}}
\toprule
方法 & SSIM 越高越好 & 调色板 W1 越低越好 & 重映射 RMSE Δ1 越低越好 & 光流 EPE Δ1 越低越好 \\
\midrule
模拟平衡方法 & 0.9395 & 8.8136 & 28.61 & 0.2963 \\
模拟闪烁方法 & 0.8270 & 7.3344 & 44.06 & 0.5494 \\
模拟覆盖方法 & 0.7122 & 8.1708 & 65.43 & 0.1389 \\
\bottomrule
\end{tabular}
\end{table}
\begin{figure}[H]
\centering
\includegraphics[width=\linewidth]{assets/demo\_case\_01.jpg}
\caption{demo\_case\_01 的盲评源视频、风格参考与候选结果比较。}
\end{figure}
\subsection*{demo\_case\_02}
\begin{table}[H]
\centering
\caption{demo\_case\_02 的代表性指标。}
\small
\begin{tabular}{@{}lcccc@{}}
\toprule
方法 & SSIM 越高越好 & 调色板 W1 越低越好 & 重映射 RMSE Δ1 越低越好 & 光流 EPE Δ1 越低越好 \\
\midrule
模拟平衡方法 & 0.9101 & 8.6103 & 39.39 & 0.2500 \\
模拟闪烁方法 & 0.7736 & 6.9334 & 59.84 & 0.8630 \\
模拟覆盖方法 & 0.7315 & 7.0425 & 69.83 & 0.2467 \\
\bottomrule
\end{tabular}
\end{table}
\begin{figure}[H]
\centering
\includegraphics[width=\linewidth]{assets/demo\_case\_02.jpg}
\caption{demo\_case\_02 的盲评源视频、风格参考与候选结果比较。}
\end{figure}
\subsection*{demo\_case\_03}
\begin{table}[H]
\centering
\caption{demo\_case\_03 的代表性指标。}
\small
\begin{tabular}{@{}lcccc@{}}
\toprule
方法 & SSIM 越高越好 & 调色板 W1 越低越好 & 重映射 RMSE Δ1 越低越好 & 光流 EPE Δ1 越低越好 \\
\midrule
模拟平衡方法 & 0.9376 & 9.6385 & 27.09 & 0.1852 \\
模拟闪烁方法 & 0.8327 & 7.9925 & 39.42 & 0.2154 \\
模拟覆盖方法 & 0.7022 & 8.9114 & 56.35 & 0.1766 \\
\bottomrule
\end{tabular}
\end{table}
\begin{figure}[H]
\centering
\includegraphics[width=\linewidth]{assets/demo\_case\_03.jpg}
\caption{demo\_case\_03 的盲评源视频、风格参考与候选结果比较。}
\end{figure}
\textit{本报告由 video-style-transfer-eval 自动生成。}
\end{document}
